\documentclass[12pt,a4paper]{article}
\usepackage[utf8]{inputenc}
\usepackage[T1]{fontenc}
\usepackage{amsmath}
\usepackage{amsfonts}
\usepackage{amssymb}
\usepackage{graphicx}
\usepackage{float}
\usepackage{hyperref}
\usepackage{natbib}
\usepackage{geometry}
\geometry{margin=1in}

\title{Uzbekistan's Economic Integration with Northeast Asia: Trade, FDI, and Global Value Chains}
\author{Your Name}
\date{\today}

\begin{document}

\maketitle

\begin{abstract}
This dissertation examines Uzbekistan's economic relations with China, Japan, and South Korea, focusing on bilateral trade flows, foreign direct investment, and integration into global value chains following the 2017 liberalization reforms. Using data from UN Comtrade, UNCTAD, and World Bank, the analysis reveals significant growth in trade and FDI, particularly with China, and explores the mechanisms through which trade openness contributes to manufacturing development and economic complexity.
\end{abstract}

\tableofcontents
\newpage

\section{Introduction}
\label{sec:intro}

Uzbekistan's 2017 economic liberalization reforms marked a significant shift towards greater integration with global markets. This dissertation investigates the impact of these reforms on relations with Northeast Asia's economic triad: China, Japan, and South Korea.

Key research questions:
\begin{itemize}
\item How have bilateral trade flows evolved post-2017?
\item What role has FDI from Northeast Asia played in Uzbekistan's development?
\item To what extent has increased trade openness facilitated integration into global value chains?
\item What are the implications for economic complexity and manufacturing development?
\end{itemize}

\section{Literature Review}
\label{sec:lit}

[Literature review content here]

\section{Methodology}
\label{sec:method}

Data sources include:
\begin{itemize}
\item Bilateral trade data from UN Comtrade
\item FDI data from UNCTAD
\item Macroeconomic indicators from World Bank WDI
\item Economic Complexity Index from OEC
\end{itemize}

Analysis employs descriptive statistics, time series analysis, and visualization techniques.

\section{Results}
\label{sec:results}

\subsection{Bilateral Trade Flows}
\label{subsec:trade}

\begin{figure}[H]
\centering
\includegraphics[width=\textwidth]{figures/fig7_bilateral_trade.pdf}
\caption{Uzbekistan Bilateral Trade with Northeast Asia Triad (2010–2023)}
\label{fig:trade}
\end{figure}

The figure shows significant growth in total trade, particularly with China, following the 2017 reforms.

\subsection{FDI Inflows}
\label{subsec:fdi}

\begin{figure}[H]
\centering
\includegraphics[width=\textwidth]{figures/fig8_fdi_by_origin.pdf}
\caption{FDI Inflows into Uzbekistan by Northeast Asian Origin (2015–2023)}
\label{fig:fdi}
\end{figure}

FDI from China dominates, with substantial increases post-reform.

\subsection{Intermediate Goods and GVC Integration}
\label{subsec:gvc}

\begin{figure}[H]
\centering
\includegraphics[width=\textwidth]{figures/fig9_intermediate_goods_share.pdf}
\caption{Share of Intermediate Goods in Uzbekistan's Imports from Northeast Asia}
\label{fig:gvc}
\end{figure}

The share of intermediate goods in imports has increased, indicating deeper GVC participation.

\subsection{Trade Openness and Economic Complexity}
\label{subsec:complexity}

\begin{figure}[H]
\centering
\includegraphics[width=\textwidth]{figures/fig10_mechanism_trajectory.pdf}
\caption{Trade Openness, Manufacturing VA, and ECI Trajectory (2010–2023)}
\label{fig:complexity}
\end{figure}

The trajectory shows improvements in trade openness and manufacturing value added, with gradual ECI gains.

\subsection{Post-Reform Trade Growth Decomposition}
\label{subsec:growth}

\begin{figure}[H]
\centering
\includegraphics[width=\textwidth]{figures/fig11_growth_decomposition.pdf}
\caption{Post-Reform Trade Growth Decomposition by Northeast Asian Partner}
\label{fig:growth}
\end{figure}

China accounts for the majority of post-reform trade growth.

\section{Discussion}
\label{sec:discussion}

[Discussion of results, implications, limitations]

\section{Conclusion}
\label{sec:conclusion}

[Conclusions and future research directions]

\bibliographystyle{plain}
\bibliography{references}

\end{document}